\documentclass[11pt,a4paper]{ctexart}
\usepackage[margin=2.2cm]{geometry}
\usepackage{amsmath,amssymb,amsthm,mathtools,booktabs,hyperref,microtype}
\hypersetup{colorlinks=true,linkcolor=black,urlcolor=blue}
\newtheorem{theorem}{定理}
\title{\bfseries 共享根深度的有限观测证书\\与采样复杂度}
\author{研究札记 · F3-REC}
\date{2026 年 9 月 29 日}
\begin{document}
\maketitle
\begin{abstract}
固定共享根深度模型
\[
R_i=v_3(d-\alpha_i),\qquad L_{ij}=v_3(\alpha_i-\alpha_j)
\]
中，本文研究只观察带标签深度向量时如何恢复根距离矩阵。核心事实极为简单：一旦某对深度不同，较小值就等于根距离，因此第一次不等观测就是无误判的精确证书。由此得到几何等待时间、完整矩阵恢复的高概率上界，以及区分距离 \(L\) 与 \(L+1\) 所需的 \(\Omega(3^L)\) 被动样本下界。进一步给出截断深度的精确恢复规则和三个等距根的停止时间闭式。状态为 \texttt{PAPER-AUDITED}，尚未 Lean 形式化。
\end{abstract}

\section{模型}
固定
\[
b\ge1,\qquad
R_i=v_3(d-\alpha_i),\qquad
D_i=b+R_i,
\]
以及
\[
L_{ij}=v_3(\alpha_i-\alpha_j).
\]
所有根是三进单位并在模 \(3\) 下同余。

局部单位域归一化 Haar 模型中，
\[
\Pr(R_i=0)=\frac12,\qquad
\Pr(R_i=t)=3^{-t}\quad(t\ge1).
\]

\section{一次不等观测就是证书}
对
\[
x=d-\alpha_i,\qquad
y=d-\alpha_j
\]
有
\[
x-y=\alpha_j-\alpha_i.
\]
强三角不等式给出
\[
\min(R_i,R_j)\le L_{ij}.
\]
若
\[
R_i\ne R_j,
\]
则不等赋值时和差取较小值，所以
\[
\boxed{
L_{ij}=\min(R_i,R_j).
}
\]

令
\[
\Delta=R_i-R_j,\qquad L=L_{ij}.
\]
若 \(\Delta=h>0\)，则
\[
(R_i,R_j)=(L+h,L),
\]
该事件概率是
\[
3^{-(L+h)}.
\]
负差同理，因此
\[
\boxed{
\Pr(\Delta=0)=1-3^{-L},
}
\]
\[
\boxed{
\Pr(\Delta=h)=3^{-(L+|h|)}\quad(h\ne0).
}
\]
所以
\[
\boxed{\Pr(R_i\ne R_j)=3^{-L_{ij}}.}
\]

\section{无误判算法与等待时间}
扫描独立观测向量。对每一对 \((i,j)\)，首次遇到
\[
R_i\ne R_j
\]
就记录
\[
L_{ij}=\min(R_i,R_j).
\]
没有不等观测时只标记“未知”。

单对等待时间 \(\tau_{ij}\) 是参数
\[
3^{-L_{ij}}
\]
的几何分布：
\[
\boxed{
\Pr(\tau_{ij}>T)
=
(1-3^{-L_{ij}})^T,
\qquad
\mathbb E\tau_{ij}=3^{L_{ij}}.
}
\]

令
\[
K=\max_{i\ne j}L_{ij},
\qquad
B_m=\binom m2.
\]
全矩阵停止时间 \(\tau\) 满足
\[
\boxed{
\Pr(\tau>T)
\le
B_me^{-T3^{-K}},
}
\]
\[
\boxed{
\mathbb E\tau
\le
1+3^K(1+\log B_m).
}
\]

\section{被动样本下界}
考虑两个可能模型，其根距分别为
\[
L,\qquad L+1.
\]
两分布总变差为
\[
\boxed{
\|\mu_L-\mu_{L+1}\|_{\rm TV}=3^{-L}.
}
\]
它们单次观测的共同质量为
\[
1-3^{-L}.
\]
所以 \(T\) 次乘积分布共同质量至少
\[
(1-3^{-L})^T.
\]
任意判别规则的两种错误概率之和至少为该共同质量，因此较大错误率至少
\[
\boxed{
\frac12(1-3^{-L})^T.
}
\]
要求两边都不超过 \(\delta<1/2\) 时，
\[
\boxed{
T
\ge
\frac{\log(1/(2\delta))}
{-\log(1-3^{-L})}
\ge
\frac23\,3^L\log(1/(2\delta)).
}
\]

\section{截断深度}
固定
\[
H\ge1,\qquad
Y_i=\min(R_i,H).
\]
目标是
\[
L^{[H]}_{ij}=\min(L_{ij},H).
\]

若
\[
Y_i\ne Y_j,
\]
则较小值未截断，
\[
L_{ij}=\min(Y_i,Y_j).
\]
若
\[
Y_i=Y_j=H,
\]
则
\[
L_{ij}\ge H.
\]
所以两种情况都给精确的截断距离证书。

单次证书概率
\[
\boxed{
\pi_H(L)=
\begin{cases}
3^{-L},&L<H,\\
1/(2\cdot3^{H-1}),&L\ge H.
\end{cases}
}
\]
于是
\[
\boxed{
\Pr(\text{\(T\) 次后未恢复全部截断距离})
\le
B_m\exp\!\left(-\frac{T}{2\cdot3^{H-1}}\right).
}
\]

当两个真实距离都至少为 \(H\) 时，截断观测分布完全相同，所以有限精度数据本来就不能恢复 \(H\) 以上的信息。

\section{三个等距根}
若
\[
m=3,\qquad
L_{12}=L_{13}=L_{23}=L,
\]
三个根占共同父球的三个下一层子球。进入两个不同子球以后，三对距离证书全部完成。

容斥得到
\[
\boxed{
\Pr(\tau>T)
=
3(1-3^{-L})^T
-
2\left(1-\frac32\,3^{-L}\right)^T,
}
\]
以及
\[
\boxed{
\mathbb E\tau
=
5\cdot3^{L-1}.
}
\]

\section{边界}
证书正确性是确定性的；只有等待时间使用独立采样模型。本文讨论被动观测，不讨论已知系数后的直接计算，也不把截断恢复冒充无限精度恢复。

\begin{thebibliography}{9}
\bibitem{Ultra} S. G. Corregidor and A. Martínez-Pérez, work on finite metric and \(k\)-metric bases in ultrametric spaces.
\end{thebibliography}
\end{document}

% final-proof-pdf-build
